\chapter{GROW Model\\โครงสร้างบทสนทนาที่นำไปสู่การลงมือทำ}

บทสนทนาที่อบอุ่นแต่ไม่เคลื่อนไปสู่การตัดสินใจอาจทำให้ทั้งสองฝ่ายรู้สึกดีชั่วคราว แต่ไม่สร้างการเติบโต GROW เป็นโครงสร้างเรียบง่ายที่ช่วยรักษาสมดุลระหว่างการสำรวจและการลงมือทำ โดยไม่บังคับให้ทุกการสนทนาเดินเป็นเส้นตรง

\learningobjectives{
  \item ใช้ Goal, Reality, Options และ Way Forward ในหนึ่ง session
  \item ปรับ GROW ให้เหมาะกับความซับซ้อนของประเด็น
  \item ปิดการสนทนาด้วยความรับผิดชอบที่ชัด
}

\section{สี่จังหวะของ GROW}
\begin{longtable}{P{.15\textwidth}P{.27\textwidth}P{.44\textwidth}}
\toprule
\textbf{จังหวะ} & \textbf{หน้าที่} & \textbf{คำถาม Anchor}\\
\midrule
Goal & กำหนดผลที่ต้องการจาก session & \enquote{เมื่อจบวันนี้ อะไรจะทำให้เวลานี้มีคุณค่า}\\
Reality & เข้าใจสถานการณ์ หลักฐาน และอุปสรรค & \enquote{ขณะนี้เกิดอะไรขึ้นจริง และเรารู้อะไรแน่}\\
Options & สร้างและประเมินทางเลือก & \enquote{มีทางเลือกใดบ้าง และแต่ละทางแลกกับอะไร}\\
Way Forward & เลือกก้าวต่อไปและความรับผิดชอบ & \enquote{ท่านจะทำอะไร ภายในเมื่อใด และต้องการการสนับสนุนใด}\\
\bottomrule
\end{longtable}

\section{บทสนทนาตัวอย่าง: Impact Statement}
\begin{examplebox}
\textbf{Mentor (Goal):} \enquote{หากจบชั่วโมงนี้พร้อมสิ่งหนึ่ง ท่านอยากได้อะไร}\\
\textbf{Mentee:} \enquote{ร่าง Impact Statement ที่คนให้ทุนเข้าใจ}\\
\textbf{Mentor (Reality):} \enquote{ฉบับปัจจุบันทำให้คนอ่านเข้าใจอะไร และยังไม่เห็นอะไร}\\
\textbf{Mentee:} \enquote{เห็นตัวเลขประสิทธิภาพ แต่ไม่เห็นว่าใครได้ประโยชน์}\\
\textbf{Mentor (Options):} \enquote{เราจะเล่าโดยเริ่มจากเกษตรกร ผู้ประกอบการ หรือนโยบาย แต่ละมุมมีหลักฐานใด}\\
\textbf{Mentee:} \enquote{หลักฐานเกษตรกรชัดที่สุด}\\
\textbf{Mentor (Way Forward):} \enquote{ท่านจะร่างสามประโยคจากมุมนั้นภายในวันไหน และใครช่วยตรวจข้อเท็จจริง}\\
\end{examplebox}

\section{GROW ไม่ใช่แบบสอบถาม}
บางครั้ง Reality ใหม่ทำให้ต้องกลับไปปรับ Goal หรือ Options หนึ่งทางเผยข้อมูลที่ต้องสำรวจเพิ่ม Mentor ควรใช้ GROW เป็นแผนที่ ไม่ใช่เส้นทางบังคับ สัญญาณสำคัญคือ Mentee กำลังคิดลึกขึ้นและเข้าใกล้การตัดสินใจที่เป็นของตน

\section{ปิดวงให้ชัด}
Way Forward ต้องเฉพาะเจาะจงพอที่จะติดตามได้: ใครทำอะไร เมื่อใด สำเร็จมีหน้าตาอย่างไร และต้องการความช่วยเหลืออะไร ก่อนจบให้ Mentee สรุปด้วยคำของตนเอง และตกลงวิธีติดตามโดยไม่ทำให้ Mentor กลายเป็นผู้ควบคุม

\begin{mentornote}
หาก Mentee ยังไม่พร้อมเลือกการกระทำ เป้าหมายของ session อาจเป็นการสร้างความเข้าใจ ไม่ควรเร่ง action เพื่อให้แบบฟอร์มครบ แต่ต้องบันทึกให้ชัดว่าสิ่งใดต้องเข้าใจก่อนตัดสินใจ
\end{mentornote}

\begin{keytakeaway}
GROW ช่วยให้บทสนทนาเคลื่อนจากความตั้งใจไปสู่ความจริง ทางเลือก และก้าวต่อไป โดยคงให้ Mentee เป็นเจ้าของเป้าหมายและการลงมือทำ
\end{keytakeaway}

\begin{reflection}
วางแผน session 45 นาทีด้วยคำถาม GROW อย่างละสองข้อ หลัง session บันทึกว่าจังหวะใดใช้เวลามากที่สุด จังหวะใดถูกรีบผ่าน และสิ่งนั้นบอกอะไรเกี่ยวกับประเด็นของ Mentee
\blankline\blankline
\end{reflection}

\chaptersummary{GROW เป็นโครงสร้างที่ยืดหยุ่นสำหรับสร้างเป้าหมาย เข้าใจความจริง เปิดทางเลือก และตกลงก้าวต่อไป การใช้ที่ดีไม่ติดกับลำดับ แต่รักษาความชัดและความเป็นเจ้าของของ Mentee}
